\section{Safety benchmarks：安全是什么也要定义}

安全评价尤其能说明本讲主旨：如果 construct 不清楚，metric 就会误导。AI safety 不是一个单指标，它包含违法伤害、隐私、偏见、幻觉、sycophancy、jailbreak、dual-use 等维度。

这一节从 crash test 类比开始，是为了强调“安全”必须场景化。汽车安全不会只给一个总分，而会区分正面碰撞、侧面碰撞、儿童乘员、行人保护。AI safety 同理：harmful request refusal、policy compliance、jailbreak robustness、dual-use capability、truthfulness 和 privacy 都是不同维度。

\begin{figure}[H]
\centering
\includegraphics[width=0.45\textwidth]{crash-test.jpg}
\caption{Crash-test analogy：安全评价像碰撞测试，需要定义风险场景和评分规则。}
\figsource{CS336 Lecture 12 官方源引用图，本地化远程图片。}
\end{figure}

\begin{knowledgebox}{读图：为什么用 crash test 类比 AI safety}
汽车安全不是一句“安全/不安全”，而是不同碰撞场景、速度、假人伤害指标。AI safety 也类似：harmful instructions、隐私、偏见、幻觉、jailbreak、dual-use 都是不同风险维度。一个模型在某类安全测试上好，不代表全局安全。
\end{knowledgebox}

\begin{figure}[H]
\centering
\includegraphics[width=0.88\textwidth]{air-overview.png}
\caption{AIR-Bench：基于监管框架和公司政策，把风险分成 314 类、5694 prompts。}
\figsource{CRFM HELM AIR-Bench 图，本地化后供 CS336 Lecture 12 使用。}
\end{figure}

\begin{knowledgebox}{读图：AIR-Bench 的 taxonomized risks}
安全评价需要 taxonomy，否则不同评测只是在测不同风险。AIR-Bench 把风险类别显式化，能帮助比较模型在不同政策维度上的拒答和合规表现。读图时要看风险分类如何组织，而不是只看总分。
\end{knowledgebox}

\begin{figure}[H]
\centering
\includegraphics[width=0.9\textwidth]{gcg-examples.png}
\caption{GCG jailbreak examples：自动优化 prompts 绕过安全拒答，并可能迁移到闭源模型。}
\figsource{CS336 Lecture 12 官方源引用图。}
\end{figure}

\begin{warningbox}{读图：jailbreak 评价是对齐系统的压力测试}
GCG 说明安全不是只看普通 harmful prompt 是否拒答，还要看 adversarial prompts 是否能绕过。迁移性意味着 open-weight model 上找到的攻击可能影响 closed models。老师在这里的提醒是：安全 eval 必须包含对抗性，而不只是正常用户输入。
\end{warningbox}

\begin{importantbox}{Safety 的上下文依赖}
很多安全问题强烈依赖政治、法律、社会规范和地区差异。比如医疗建议、网络安全、政治内容、犯罪协助在不同上下文中的边界不同。安全 benchmark 必须明确政策背景，而不是假装存在唯一全局答案。
\end{importantbox}

